\chapter{Thiết lập thử nghiệm}
    \section{Bộ dữ liệu}
        Các thử nghiệm được thực hiện trên ba bộ dữ liệu phân loại ảnh tiêu chuẩn.
        Các bộ dữ liệu này được chọn vì chúng được sử dụng rộng rãi trong văn
        献 nghiên cứu về pruning mạng nơ-ron, cho phép so sánh trực tiếp với
        các công trình trước đây \cite{imp, grasp, synflow, hybrid}.
         
        \subsection{MNIST}
        \label{subsec:mnist}
         
        MNIST~\cite{lecun1998mnist} là bộ dữ liệu nhận dạng chữ số viết tay gồm 70\,000
        ảnh thang độ xám kích thước $28 \times 28$ pixel, trong đó 60\,000 ảnh dùng để
        huấn luyện và 10\,000 ảnh dùng để kiểm tra.
        Bộ dữ liệu bao gồm 10 lớp tương ứng với các chữ số từ 0 đến 9.
         
        Trong luận văn này, MNIST được sử dụng như một thiết lập kiểm chứng nhanh
        (\textit{quick validation}) để xác nhận tính đúng đắn của các thuật toán pruning
        trước khi chạy trên các bộ dữ liệu phức tạp hơn.
        Bộ dữ liệu không yêu cầu tiền xử lý đặc biệt; ảnh chỉ được chuẩn hóa về
        khoảng $[0, 1]$ bằng cách chia cho 255.
         
        % ----------------------------------------------------------
        \subsection{CIFAR-10}
        \label{subsec:cifar10}
         
        CIFAR-10~\cite{krizhevsky2009learning} là bộ dữ liệu phân loại ảnh màu
        kích thước $32 \times 32$ pixel gồm 60\,000 ảnh thuộc 10 lớp đối tượng
        (máy bay, ô tô, chim, mèo, hươu, chó, ếch, ngựa, tàu thủy, xe tải),
        với 50\,000 ảnh huấn luyện và 10\,000 ảnh kiểm tra (mỗi lớp 5\,000/1\,000 ảnh).
         
        CIFAR-10 là bộ dữ liệu tham chiếu chính trong luận văn vì đây là benchmark
        tiêu chuẩn được hầu hết các công trình \gls{imp} sử dụng.
        
        Quá trình tiền xử lý bao gồm:
        \begin{itemize}
            \item \textbf{Huấn luyện:} Lật ngang ngẫu nhiên (random horizontal flip),
                  cắt ngẫu nhiên sau padding 4 pixel (random crop with padding=4),
                  chuẩn hóa với giá trị trung bình $\mu = (0.4914,\,0.4822,\,0.4465)$
                  và độ lệch chuẩn $\sigma = (0.2023,\,0.1994,\,0.2010)$.
            \item \textbf{Kiểm tra:} Chỉ chuẩn hóa (không tăng cường dữ liệu).
        \end{itemize}
         
        % ----------------------------------------------------------
        \subsection{CIFAR-100}
        \label{subsec:cifar100}
         
        CIFAR-100~\cite{krizhevsky2009learning} có cùng định dạng và kích thước với
        CIFAR-10 (60\,000 ảnh $32 \times 32$) nhưng được phân chia thành 100 lớp chi
        tiết hơn, mỗi lớp có 500 ảnh huấn luyện và 100 ảnh kiểm tra.
        Đây là bộ dữ liệu thách thức hơn vì không gian nhãn rộng hơn đáng kể.
         
        CIFAR-100 được sử dụng để đánh giá khả năng mở rộng (scalability) của
        các phương pháp pruning khi độ phức tạp của tác vụ phân loại tăng lên.
        Quy trình tiền xử lý tương tự CIFAR-10, với giá trị chuẩn hóa riêng:
        $\mu = (0.5071,\,0.4867,\,0.4408)$ và
        $\sigma = (0.2675,\,0.2565,\,0.2761)$.
         
        \begin{table}[htbp]
        \centering
        \caption{Tổng quan các bộ dữ liệu sử dụng trong thực nghiệm.}
        \label{tab:datasets}
        \begin{tabular}{lcccccc}
        \hline
        \textbf{Bộ dữ liệu} & \textbf{Lớp} & \textbf{Kích thước ảnh} & \textbf{Train} & \textbf{Test} & \textbf{Màu sắc} & \textbf{Vai trò} \\
        \hline
        MNIST     & 10  & $28\times28$ & 60\,000 & 10\,000 & Xám & Kiểm chứng nhanh \\
        CIFAR-10  & 10  & $32\times32$ & 50\,000 & 10\,000 & RGB & Benchmark chính  \\
        CIFAR-100 & 100 & $32\times32$ & 50\,000 & 10\,000 & RGB & Đánh giá mở rộng \\
        \hline
        \end{tabular}
        \end{table}

    \section{Kiến trúc mạng}
    Bốn kiến trúc mạng nơ-ron được sử dụng trong luận văn nhằm đánh giá
hiệu quả của các thuật toán pruning trên nhiều quy mô mô hình khác nhau,
từ mạng nhỏ gọn đến mạng sâu có dung lượng lớn.
 
    % ----------------------------------------------------------
    \subsection{ResNet-20}
    \label{subsec:resnet20}
     
    ResNet-20~\cite{he2016deep} là phiên bản nhỏ nhất của họ ResNet được thiết kế
    đặc biệt cho bộ dữ liệu CIFAR.
    Mạng gồm 20 lớp tích chập với tổng cộng khoảng \textbf{270\,000 tham số},
    được tổ chức thành ba nhóm (stages) với lần lượt 16, 32 và 64 bộ lọc,
    mỗi nhóm gồm $n=3$ khối residual.
     
    Cơ chế kết nối tắt (skip connection) trong mỗi khối residual cho phép
    gradient lan truyền hiệu quả, giúp mạng hội tụ ổn định ngay cả sau khi
    cắt tỉa lớn.
    ResNet-20 đóng vai trò kiến trúc tham chiếu chính (primary baseline) và
    được sử dụng trong tất cả các thử nghiệm so sánh do chi phí tính toán thấp,
    phù hợp với việc chạy nhiều cấu hình thử nghiệm.
     
    % ----------------------------------------------------------
    \subsection{ResNet-50}
    \label{subsec:resnet50}
     
    ResNet-50~\cite{he2016deep} là phiên bản sâu hơn với 50 lớp, sử dụng
    kiến trúc khối \textit{bottleneck} (conv $1\times1$ -- conv $3\times3$ -- conv $1\times1$)
    với khoảng \textbf{25,6 triệu tham số}.
    Mạng gồm bốn stages với lần lượt 64, 128, 256 và 512 kênh đặc trưng
    (feature channels).
     
    ResNet-50 được sử dụng trên CIFAR-100 để đánh giá liệu các \textit{winning
    ticket} có tồn tại và duy trì hiệu năng cao trong mô hình có dung lượng lớn
    hay không.
    Khối lượng tham số cao hơn cũng giúp kiểm tra khả năng mở rộng của
    thuật toán Genetic Algorithm và Hybrid-Improve khi không gian tìm kiếm tăng lên.

    \section{Siêu tham số}
    \label{subsec:algo_hparams}
        
        \begin{table}[htbp]
        \centering
        \caption{Siêu tham số đặc thù cho từng thuật toán pruning (CIFAR-10/ResNet-20).}
        \label{tab:algo_hparams}
        \resizebox{\textwidth}{!}{%
        \begin{tabular}{llll}
        \hline
        \textbf{Thuật toán} & \textbf{Siêu tham số} & \textbf{Giá trị} & \textbf{Ý nghĩa} \\
        \hline
        \multirow{7}{*}{IMP}
        & \texttt{target\_sparsity}     & 0.3, 0.6          & Tỷ lệ tham số bị cắt tỉa \\
        & \texttt{num\_iterations}      & 5                 & Số vòng lặp pruning \\
        & \texttt{epochs\_per\_iteration} & 100              & Số epoch mỗi vòng lặp \\
        & \texttt{use\_global\_pruning} & True              & Pruning toàn cục hay theo lớp \\
        & \texttt{learning\_rate}       & 0.1               & Tốc độ học \\
        & \texttt{batch\_size}          & 128               & Kích thước batch \\
        & \texttt{weight\_decay}        & 0.0005            & Hệ số regularization \\
        \hline
        \multirow{5}{*}{Early-Bird}
        & \texttt{target\_sparsity}     & 0.3, 0.6          & Tỷ lệ kênh bị cắt tỉa \\
        & \texttt{search\_epochs}       & 17--19            & Số epoch tìm kiếm vé (thực tế) \\
        & \texttt{finetune\_epochs}     & 20--100+          & Số epoch tinh chỉnh (thực tế) \\
        & \texttt{l1\_coef}             & $10^{-4}$         & Hệ số phạt L1 trên $\gamma$ BN \\
        & \texttt{distance\_threshold}  & 0.1               & Ngưỡng hội tụ mặt nạ \\
        \hline
        \multirow{8}{*}{GraSP}
        & \texttt{target\_sparsity}     & 0.3, 0.6          & Tỷ lệ tham số bị cắt tỉa \\
        & \texttt{epochs}               & 160               & Số epoch huấn luyện sau pruning \\
        & \texttt{samples\_per\_class}  & 10                & Số mẫu/lớp để tính GraSP score \\
        & \texttt{grasp\_T}             & 200.0             & Nhiệt độ scaling logits \\
        & \texttt{grasp\_iters}         & 1                 & Số lần tích lũy gradient \\
        & \texttt{learning\_rate}       & 0.1               & Tốc độ học \\
        & \texttt{batch\_size}          & 128               & Kích thước batch \\
        & \texttt{weight\_decay}        & 0.0001            & Hệ số regularization \\
        \hline
        \multirow{8}{*}{SynFlow}
        & \texttt{rho}                  & 1.43, 2.5         & Tỷ lệ nén $\rho$ (tương ứng sparsity 0.3, 0.6) \\
        & \texttt{target\_sparsity}     & 0.3, 0.6          & Tỷ lệ tham số bị cắt tỉa (thực tế) \\
        & \texttt{synflow\_iters}       & 100               & Số vòng lặp SynFlow \\
        & \texttt{epochs}               & 160               & Số epoch huấn luyện sau pruning \\
        & \texttt{learning\_rate}       & 0.1               & Tốc độ học \\
        & \texttt{batch\_size}          & 128               & Kích thước batch \\
        & \texttt{weight\_decay}        & 0.0001            & Hệ số regularization \\
        \hline
        \multirow{11}{*}{Hybrid}
        & \texttt{target\_sparsity}     & 0.3, 0.6          & Tỷ lệ cắt tỉa mục tiêu \\
        & \texttt{oneshot\_ratio}       & 0.7               & Phần one-shot của mục tiêu \\
        & \texttt{iterative\_step}      & 0.1               & Phần trọng số cắt mỗi bước lặp \\
        & \texttt{initial\_epochs}      & 160               & Epoch huấn luyện dày đặc \\
        & \texttt{initial\_lr}          & 0.01              & LR ban đầu pha huấn luyện dày đặc \\
        & \texttt{oneshot\_finetune\_max\_epochs} & 200 & Epoch tinh chỉnh sau one-shot \\
        & \texttt{oneshot\_finetune\_patience} & 200 & Patience tinh chỉnh one-shot \\
        & \texttt{iter\_finetune\_max\_epochs}  & 10  & Epoch tinh chỉnh mỗi bước lặp \\
        & \texttt{iter\_finetune\_patience} & 10     & Early stopping bước lặp \\
        & \texttt{batch\_size}          & 128               & Kích thước batch \\
        & \texttt{weight\_decay}        & 0.0005            & Hệ số regularization \\
        \hline
        \multirow{15}{*}{Hybrid-Improve}
        & \texttt{target\_sparsity}     & 0.3, 0.6          & Tỷ lệ cắt tỉa mục tiêu \\
        & \texttt{oneshot\_ratio}       & 0.79              & Phần one-shot thích ứng \\
        & \texttt{iterative\_step}      & 0.2               & Phần trọng số cắt mỗi bước lặp \\
        & \texttt{initial\_epochs}      & 160               & Epoch huấn luyện dày đặc \\
        & \texttt{initial\_lr}          & 0.1               & LR ban đầu pha huấn luyện dày đặc \\
        & \texttt{oneshot\_finetune\_max\_epochs} & 200 & Epoch tinh chỉnh sau one-shot \\
        & \texttt{oneshot\_finetune\_patience} & 200 & Patience tinh chỉnh one-shot \\
        & \texttt{iter\_finetune\_max\_epochs}  & 10  & Epoch tinh chỉnh mỗi bước lặp \\
        & \texttt{iter\_finetune\_patience} & 10     & Early stopping bước lặp \\
        & \texttt{batch\_size}          & 512               & Kích thước batch \\
        & \texttt{weight\_decay}        & 0.0005            & Hệ số regularization \\
        & \texttt{use\_fisher\_adaptive\_ratio} & True & Sử dụng Fisher adaptive ratio \\
        & \texttt{use\_kd}              & True              & Sử dụng knowledge distillation \\
        & \texttt{kd\_temperature}       & 4.0               & Nhiệt độ KD \\
        & \texttt{kd\_lambda}           & 0.5               & Hệ số KD \\
        \hline
        \end{tabular}%
        }
        \end{table}

    \section{Chỉ số đánh giá}
    \label{sec:metrics}
 
    Để so sánh toàn diện giữa các thuật toán pruning, luận văn sử dụng
    ba nhóm chỉ số đánh giá: hiệu năng dự đoán, mức độ nén mô hình,
    và hiệu quả tính toán.
 
% ----------------------------------------------------------
    \subsection{Độ Chính Xác Phân Loại}
    \label{subsec:accuracy}
    
    \textbf{Top-1 Test Accuracy} là chỉ số hiệu năng chính, đo tỷ lệ phần trăm
    mẫu kiểm tra được phân loại đúng:
    \begin{equation}
        \text{Acc} = \frac{\text{Số mẫu dự đoán đúng}}{\text{Tổng số mẫu kiểm tra}} \times 100\%
    \end{equation}
    
    Để đảm bảo kết quả đáng tin cậy, mỗi cấu hình thử nghiệm được chạy với
    nhiều random seed và kết quả được báo cáo dưới dạng giá trị trung bình
    $\pm$ độ lệch chuẩn.
    
    \textbf{Accuracy Drop} ($\Delta$Acc) đo mức độ suy giảm độ chính xác so
    với mô hình gốc (dense network) tại cùng kiến trúc và bộ dữ liệu:
    \begin{equation}
        \Delta\text{Acc} = \text{Acc}_{\text{dense}} - \text{Acc}_{\text{pruned}}
    \end{equation}
    Giá trị $\Delta$Acc nhỏ cho thấy phương pháp pruning bảo toàn tốt
    hiệu năng gốc.
    
    % ----------------------------------------------------------
    \subsection{Chỉ Số Nén Mô Hình}
    \label{subsec:compression_metrics}
    
    \textbf{Sparsity} ($s$) là tỷ lệ trọng số bị cắt tỉa (đặt về 0) trên
    tổng số trọng số của mô hình:
    \begin{equation}
        s = \frac{|\{w_i : w_i = 0\}|}{|\{w_i\}|} \times 100\%
    \end{equation}
    Sparsity 90\% nghĩa là chỉ còn 10\% trọng số không bằng 0.
    
    \textbf{Compression Ratio} ($\rho$) biểu diễn mức độ nén theo tỷ lệ:
    \begin{equation}
        \rho = \frac{1}{1 - s}
    \end{equation}
    Ví dụ, sparsity 90\% tương đương $\rho = 10\times$ (mô hình được nén 10 lần).
    
    \textbf{Parameters Remaining} ($P_r$) là số lượng tuyệt đối hoặc tỷ lệ
    trọng số còn lại sau pruning:
    \begin{equation}
        P_r = (1 - s) \times P_{\text{total}}
    \end{equation}
